\documentclass{kuisthesis}

\usepackage[dvipdfmx]{graphicx}
\usepackage{float}

\jtitle{粘土による直感的操作と生成AIを併用したデザイン探索手法}
\etitle{Design Exploration Method Combining Intuitive Clay Manipulation with Generative AI}
\jauthor{梅木敦成}
\eauthor{Taisei Umeki}
\supervisor{中村裕一 教授}
\date{令和8年2月2日}
\begin{document}
\maketitle
\begin{jabstract}
本研究は、非専門家が持つ曖昧なデザインイメージを具体化するために、手元の粘土の操作と生成AIを組み合わせた探索的デザイン支援手法を提案する。 
従来のテキスト指示による画像生成では、言葉にしにくい微妙な形状のニュアンスを伝えることが難しいという課題がある。 
加えて、ユーザーが用いる言葉の多義性がAIとの解釈のズレを生む要因となっており、ユーザー自身がその曖昧さを自覚していないため意図通りの生成結果が得られにくいという問題も存在する。
提案手法では、(1)粘土模型の写真とテキストを組み合わせた入力、(2)AIによる「解釈の提案」を通じた言葉の多義性の解消、(3)属性値の調整による探索的な画像生成により、これらの課題の解決を図った。 
被験者2名による定性的評価の結果、粘土で手を動かすことがアイデアの発想を助ける点や、AIからの提案が探索の足がかりとして有効であることが示唆された。 
一方で、詳細な修正を行いたい場面での制御性の低さや、ユーザーの造形スキル不足により意図した形状を正確に作成できず、その不完全な形状特徴が生成結果の悪化に直結してしまう課題も確認された。 
本研究は、AIとの対話を通じて、人間の身体的な直感とAIの表現力を共存させる新たな創作プロセスを提示するものである。
\end{jabstract}
\begin{eabstract}
This study proposes an exploratory design support method that combines clay-based shape input with generative AI to help non-expert users materialize vague design ideas. 
Conventional text-to-image generation methods face challenges in conveying subtle geometric nuances solely through natural language. 
Furthermore, semantic ambiguity in user prompts often leads to a mismatch between the AI's interpretation and the user's intent, a discrepancy users are frequently unaware of, which hinders the generation of desired results. 
To address these issues, our method integrates (1) multimodal input using clay photos and text, (2) resolution of linguistic ambiguity through "interpretation suggestions" by AI, and (3) exploratory image generation based on attribute constraints. 
Qualitative evaluation with two participants suggested that hands-on clay manipulation supports ideation and that AI suggestions serve as an effective starting point for exploration. 
However, challenges were identified regarding low controllability in fine-tuning specific details and the difficulty for users to accurately mold their intended shapes due to lack of modeling skills, which directly led to unintended generation results. 
This research presents a novel creative process that harmonizes human physical intuition with AI's expressive power through interactive dialogue.
\end{eabstract}
\tableofcontents %目次の出力

% 1章
\section{はじめに}
ものづくりやプロダクトデザインにおける創作活動では、制作の初期段階において「なんとなくこういうものが作りたい」という曖昧なイメージが出発点となることが多い。
しかし専門的な訓練を受けていない非専門家にとって、頭の中のイメージを具体化し、外部へ表現することは容易ではない。
従来はスケッチやCADツールが用いられてきたが、これらは描画能力や操作スキルを要求するため、非専門家にとって難しいプロセスとなっていた。

近年、Text-to-Image を代表とする生成AIの発展により、言語（プロンプト）を入力として多様な画像を生成できるようになった。
生成AIは「描く技術」を担うことが期待されており、アイデアの視覚化を加速させている。
つまり、言語(プロンプト)インタフェースによって、描画や造形のプロセスを生成AIに任せられるようになるのではないかと期待されている。

しかし、生成AIを利用してイメージを具現化するためには、まだ多くの課題が残っている。
本研究ではその課題を以下の二点に整理する。

第一に、言語を用いて「具体的な形状」を伝えることの困難さである。言語は「雰囲気（画風やトーン）」や「概念（物体名や比喩など）」を伝えることには適しているが、「どこをどのように」曲げたいか「どの程度まで」ボリュームを増やしたいかといった空間的な特徴を正確に指示することには困難が伴う。
その結果、ユーザーは言語化できない形状へのこだわりを妥協するか、あるいは偶然生成される出力に依存して試行錯誤を繰り返す状況に陥りやすい。
これは、連続的な『形の空間』と離散的な『言葉の空間』が1対1に対応していない（言葉では表現しきれない形が無数に存在する）ことを意味する。

第二に、言語そのものが持つ「多義性（複数の異なる意味に解釈できる性質）」の問題である。
例えば単に「かっこいい」と指示しても、それが「未来的なかっこよさ」なのか「伝統的なかっこよさ」なのかは人によって異なる。 
重要な点は、ユーザー自身もこの「言葉のあいまいさ」に自覚的でない場合が多いことである。 
自分が入力した言葉が多様な解釈を含んでいることに気づかないまま生成を行うため、AIが出力する結果に対して「何かが違う」と感じつつも、どう修正すればよいか分からなくなってしまうのである。

こうした「形状記述の困難さ」と「意味解釈の多義性」により、ユーザーが頭の中で描くデザインと、AIが生成する出力との間にはどうしてもズレが生じやすく、納得のいく結果を得るまでのプロセスが長期化する傾向にある。
% これは、連続的な『形の空間』と離散的な『言葉の空間』が1対1に対応していない（言葉では表現しきれない形が無数に存在する）ことを意味する。

本研究の目標は、非専門家ユーザーの曖昧で多義的な初期イメージを、システムとの対話を通じて段階的に言語化・具体化し、最終的に明確な要求仕様へと昇華させることである。
そのために、言語化しにくい形状特徴を直接入力できる物理的インタフェースとして粘土模型を用い、その画像とテキストを併用したマルチモーダル入力を採用する。
さらに、LLMによる「解釈の提案」を行うことで、ユーザーに言葉の多義性に気づかせ、曖昧な入力テキストの意味を明確化する枠組みを設計する。
これにより、生成された多様な候補の中からユーザーの意図に沿うものを段階的に絞り込み、納得感のあるデザイン探索を可能にする。

本研究には二つの社会的意義がある。
第一に、高度な造形スキルを持たない人々が自身の内面にあるイメージを正確に外部化できるよう支援することで、個人の創造的な表現活動を促進する点である。
第二に、AIによる自動化が進む現代において、人間が本来持つ「手で考える（触覚的思考）」能力を排除するのではなく、AIによる解釈支援と組み合わせることで、人間の身体的直感とAIの表現力を共存させる新たな創作プロセスを提示する点である。

本論文の構成を述べる。
第2章では、生成AIを用いた物品デザインの背景と、探索的デザイン支援における課題を整理する。
第3章では、提案するデザインプロセスを段階的に説明し、物理入力とデジタル探索を往復する反復ループを示す。
第4章では、ユーザー実験に基づき、提案手法の有用性と課題を定性的に分析する。
第5章では、結論と今後の展望を述べる。



% 2章
\section{生成AIを用いた物品のデザイン}
\subsection{デザインにおける生成AIの利用}
% デザインプロセスにおける「探索」の重要性
デザインとは、完成された設計図を最初から出力する作業ではなく、試行錯誤を通じて頭の中の不明瞭なイメージを徐々に明確化していく探索的プロセスである。
特に初期段階では、ユーザー自身も最終的な正解像を明確に持っていないことが多く、生成AIを用いるデザインでは、AIからのフィードバックを見ながら意図を更新していく対話的なプロセスが重要となる。

% 言語インタフェースの問題
現在主流のText-to-Image手法は自然言語を入力とし、ユーザーはプロンプトを通じて生成結果を制御する。
言語は「雰囲気」や「概念」といった抽象的情報を伝えることができる一方で、「具体的な形状」や「空間的な構成」といった非言語的情報を正確に指示することは難しい。
例えば、「右側をなめらかに隆起させる」といった局所的な形状操作をテキストだけで一意に伝えることは困難であり、ユーザーが意図を持っていても出力が不安定になる。
この「伝えたい形状イメージ」と「入力手段としての言語」の間のギャップが、意図通りの生成を妨げる要因となる。

形状以外の属性についてもユーザーが頭の中に描いているイメージを正確にシステムに伝えることは簡単ではない。
つまり、具体的な意図を持っているにもかかわらず、多義的な言葉を用いてしまうと、その意図が正確に伝わらないことがある。
例えば、「かっこいい」車が示唆する具体的な属性は、ユーザーによって大きく異なり、「角ばった」、「車高の低い」、「黒い色調の」など、種々のものが考えられる。
本質的な問題は、ユーザー自身がこの言葉の多義性に無自覚である場合にある。
ユーザーは自身の言葉が意図通りに伝わると信じて入力を行うため、AIの解釈が自身の期待するデザインとずれた場合、なぜ意図が伝わらないのかを理解できず、適切な修正を行うことが困難となる。

以上より、生成AIを用いたデザインのための探索的なプロセスをサポートするメカニズムが必要である。
さらに、ユーザーに自身の入力テキストの曖昧さを「意識」させた上でユーザーの意図する特徴を先鋭化し、具体的にシステムへ伝達するためのインタフェースが必要不可欠である。

\subsection{探索的デザイン支援システム}
本研究では、物理的な形状入力と対話的な探索を融合したデザイン支援システムを提案する。
本システムは、粘土による形状入力、LLM(Large Language Model)による多義性解消（解釈提案）、および制約に基づく画像生成探索の三段階で構成される。
これにより、形状と言語の両面からユーザーの頭の中のイメージを反映させたデザインの実現を目指す。

このために以下のハイブリッドなアプローチをとる。

第一に、物理的インタフェース（粘土）による非言語的な意図伝達を行う。
言語化できない形状のニュアンス（曲がり具合、ボリューム感など）を、粘土という可塑性のある物理メディアを通じて直接入力することで、言語プロンプトの限界を補完する。
第二に、構造的解釈によって多義性を解消する。ユーザーの入力テキストに含まれる曖昧さを、AIが一義的な語彙を用いて再定義し、複数の解釈案を提示する。
ユーザーがそれを選択することで方向性の合意形成を行い、不明瞭なイメージの外在化を促す。
第三に、制約に基づく探索的ナビゲーションにより、広大なデザイン空間の中からユーザーの好む領域を探索的に絞り込む。
一発で正解を出そうとするのではなく、「もっと鋭く」「もっと有機的に」といった相対的な制約を与えることで探索を継続可能にする。

この探索プロセスを実現するために、本研究では検索支援システムFindMeで提案された『制約に基づくナビゲーション』の概念を応用し、生成モデルの探索へと拡張する。
また、意図した範囲内（制約内）で多様な候補を提示することで、ユーザー自身も気づいていなかった好みの発見（発想の拡張）を促すことを狙う。

関連研究として、StyleCLIPなど、テキスト入力によって画像の特徴を操作する手法が提案されている。
しかし、これらはテキストの意味内容との整合性（意味的な整合性）を重視しており、ユーザー個人の期待するデザインとの整合性や主観的満足度を直接支援する枠組みとは目的が異なる。
本研究は、非専門家が曖昧なイメージを段階的に具体化していくプロセスに焦点を当て、その過程に必要な入力・解釈・探索の支援を統合する点に特徴がある。



% 3章
\section{デザインプロセスの設計}
本章では、提案手法を用いた具体的なデザインプロセスの設計について述べる。
本プロセスは、ユーザーの曖昧なイメージを段階的に明確化する「意図の詳細化」と、その意図の範囲内で候補を生成しながら満足できる解を探る「デザイン空間内の探索」を往復させることで進行する。
生成AIは多数の候補を提示できる一方、提示が単なるランダムな変形の連続になってしまうと、ユーザーは「どの方向に進んでいるのか」を見失い、試行錯誤のコストが増大する。
そこで本研究では、入力の構造化と解釈の合意形成を導入し、探索の方向性を段階的に固定しながら詳細化できるように設計する。

\subsection{デザインプロセスの流れ}
本システムにおけるデザイン探索は、図\ref{fig:system-flow-entire}に示すように、Phase A〜Dを基本単位として反復する。
ユーザーは編集対象の粒度に応じて「全体編集」「部分編集」「合成」のいずれかを選択する。
全体編集は物体全体の印象や構成を扱い、部分編集は特定部位の局所的な変更を扱い、合成は部品と本体の接合・配置を扱う。

まず Phase A では、ユーザーは粘土模型の画像と、デザイン意図を表すテキストを入力する（図\ref{fig:system-flow-a}）。
次に Phase B では、入力されたクエリの多義性を解消するために、デザインや接合方法に関する複数の解釈案を提示し、ユーザーが選択することで方向性の合意形成を行う（図\ref{fig:system-flow-bc}）。
Phase C では、選択された解釈を、画像生成エンジンが処理可能な数値表現（特徴ベクトル）へ変換し、候補画像を生成して提示する。
Phase D では、提示された画像を見ながらユーザーが特定の属性に対する修正指示（制約）を追加し、その指示を満たす範囲でパラメータ（特徴ベクトル）を更新して再生成を行う（図\ref{fig:system-flow-d}）。
これにより、ユーザーは制約による『意図の反映』と、AIの生成による『思いがけない発見』を繰り返しながら、意図に近いデザインへと候補を絞り込んでいく。

また、本システムはデジタル上での完結を前提としない。デジタル上で得られた生成画像は、物理模型を修正するための参照としても利用される。ユーザーは生成画像を手がかりに粘土模型を作り直し、再び Phase A へ入力することで、物理とデジタルを往復する探索を行う。

\subsection{デザイン意図の詳細化}
デザインの初期段階においてユーザーのイメージは曖昧であり、そのままではシステムが処理可能な指示になりにくい。本節では、入力を構造化し、意図の方向性を合意形成するまでのプロセスを述べる。

Phase A（マルチモーダル入力と構造化）では、粘土模型の写真とテキストを同時に扱うことで、言語だけでは表現しにくい形状のニュアンスを保持する。
全体編集では粘土模型の画像全体を入力として扱い、物体全体の形状特徴や構成をベースラインとして認識する。
部分編集では、画像に加えて編集したい箇所をマスクとして指定し、当該領域に対する局所的な変更指示を受け付ける。
合成では、部品と本体が写った画像を入力し、それらの接合方法や配置に関する指示を受け付ける。
この3モードを切り替えることで、「全体の印象を変える」「特定箇所だけ直す」「部品を組み合わせる」といった粒度の異なる要求を同一の枠組みで扱う。

\begin{table}[H]
  \centering
  \caption{編集モードと入力の対応}
  \label{tab:edit-modes}
  \begin{tabular}{|p{0.18\linewidth}|p{0.34\linewidth}|p{0.40\linewidth}|}
    \hline
    モード & ユーザー入力 & 主な目的・適用場面 \\
    \hline
    全体編集 & 粘土画像＋テキスト & 物体全体の印象を探索し、方向性を決める。 \\
    \hline
    部分編集 & 粘土画像＋マスク＋テキスト & 特定部位の変更（取っ手、脚部など）を局所的に反映し、微調整に用いる。 \\
    \hline
    合成 & 部品と本体が写った画像＋テキスト & パーツの接合方法や配置を探索し、部品の追加・組み換えを支援する。 \\
    \hline
  \end{tabular}
\end{table}

入力テキストに固有名詞（例：「チューリップのように」）が含まれる場合、本研究ではこれを「モチーフ」として抽出し、属性情報とは分離して管理する。
固有名詞は具体的な形状イメージ（What）を指す可能性が高く、「和風」「金属的」等の抽象的な様式特徴（How）とは異なる情報であるためである。
モチーフを独立に扱うことで、ユーザーが参照したい対象を保持しつつ、属性による制御と衝突しにくい入力表現を得ることを狙う。

Phase B（意図の解釈と提案）では、非専門家が入力する抽象的なクエリの多義性を解消するために、システムが複数の解釈案を提示する。
例えば「和風にして」という指示は、落ち着いた木質の和風から、祝祭的で装飾的な和風まで幅広い解釈が可能である。
そこでシステムは、粘土画像を視覚的コンテキストとして利用しつつ、クエリに関連する代表的な属性の組み合わせを複数パターン生成して提示する。
ユーザーはその中から自身の頭の中のイメージに最も近いものを選択し、方向性の合意形成を行う。
ここで「属性」とは「金属的」「角が丸い」等の視覚的特徴であり、属性値はその適合度合いを0から1の実数値で表したものである。

\subsection{デザイン意図に基づく探索}
意図の方向性が定まった後も、細部の調整においては言語化が難しい微細な意図が残る。
本節では、属性空間内での局所的な探索によってデザインを洗練させるプロセスを述べる。

Phase C（特徴ベクトル生成と画像生成）では、Phase Bで合意された解釈（自然言語）を、画像生成エンジンが制御可能な数値表現である特徴ベクトルへ変換する。
特徴ベクトルは各属性の強度を連続値で持つ。特徴ベクトルおよび抽出されたモチーフ情報をプロンプトに組み込み、画像生成エンジンを用いて具体的なデザイン案を提示する。
編集モードに応じて、画像全体の特徴を表すグローバルベクトル（全体編集・合成で使用）と、特定部位の特徴を表す部分ベクトル（部分編集で使用）を使い分ける。
このとき部分ベクトルはマスクされた領域に対してのみ適用され、周辺のコンテキストを保持しながら局所変形を行うことを想定する。

Phase D（制約探索によるデザイン洗練）では、ユーザーは生成画像を見ながら、特定の属性に対して制約を追加することで探索空間を絞り込む。
制約は、例えば「もっと鋭利にしたい」という要求を「鋭利さ（Sharpness）の属性値が $P \ge 0.8$ を満たす」といった不等号で表現する。
システムは現在の特徴ベクトル周辺を探索しつつ、追加された制約条件を満たす部分空間を特定し、その中で特徴ベクトルを更新して画像を再生成する。
これにより、意図を満たしながらも多様なバリエーションを提示し、ユーザー自身も想定していなかった好みの発見や発想の拡張を促す。

探索を繰り返す過程では、同じ方向の変化が連続してしまい、候補が局所的に収束して多様性が失われることがある。
そこで本研究では、過去にユーザーが「好ましくない」として棄却した候補の周辺を避けるために、棄却済み候補から離れる方向への補正（斥力）を導入する。
具体的には、棄却された候補の特徴ベクトルとの差分を用い、更新時にそれらから離れる方向へベクトルをシフトさせることで、似通った失敗作の再生成を抑制する。
さらに、各生成結果をノードとして保持し、ユーザーが任意の段階に戻って再探索できるようにすることで、探索が行き詰まった場合のリカバリを容易にする。

\subsection{詳細化と探索の繰り返しによるデザイン}
以上のプロセスにより、ユーザーはデジタル空間での探索ループを回すことができる。
生成画像に対して、全体的な印象や方向性を刷新したい場合は Phase A に戻ってクエリや解釈を再生成し、形状のバリエーションや細部のニュアンスを探索したい場合は Phase D での制約追加を繰り返す。

さらに、本システムは物理とデジタルの相互作用を含む。
デジタル上で得られた生成画像が、ユーザーにとって「この方向が良い」「この形は避けたい」といった判断材料となり、粘土模型の修正へとフィードバックされる。
例えば椅子のデザインでは、初期の粘土模型で全体の包まれ感を表現し、生成画像を見ながら背もたれの厚みや脚部の広がりを粘土で調整する。
修正後の模型を再度入力することで、システムは更新された形状を初期形状として認識し、よりユーザーの意図を正確に反映した提案が可能となる。
この「物理入力 $\to$ AIによる拡張 $\to$ 物理修正 $\to$ 再入力」というサイクルの反復により、曖昧な初期イメージは実装可能なレベルの具体的デザインへと段階的に昇華される。

\begin{figure}[H]
  \centering
  \includegraphics[width=0.65\linewidth,keepaspectratio]{/Users/umeking5/university/lab_thesis/flows/flow_entire.png}
  \caption{本システムの全体フロー}
  \label{fig:system-flow-entire}
\end{figure}

\begin{figure}[H]
  \centering
  \includegraphics[width=\textwidth]{/Users/umeking5/university/lab_thesis/flows/flow_a.drawio.png}
  \caption{フェーズA：入力}
  \label{fig:system-flow-a}
\end{figure}

\begin{figure}[H]
  \centering
  \includegraphics[width=\textwidth]{/Users/umeking5/university/lab_thesis/flows/flow_bc.drawio.png}
  \caption{フェーズB：解釈提案 / フェーズC：特徴ベクトル生成と画像生成}
  \label{fig:system-flow-bc}
\end{figure}

\begin{figure}[H]
  \centering
  \includegraphics[width=\textwidth]{/Users/umeking5/university/lab_thesis/flows/flow_d.drawio.png}
  \caption{フェーズD：制約探索によるデザイン洗練}
  \label{fig:system-flow-d}
\end{figure}




% 4章
\section{実験例と考察}
本章では、提案手法（粘土による形状入力と生成AIを活用した探索）が、デザインの専門知識を持たないユーザーにとってデザインを得る助けになっていることを検証するための実験について述べる。
比較のため、(1)粘土＋テキスト入力の後にテキストプロンプトのみで反復生成する手法（比較A）、(2)テキストのみで反復生成する手法（比較B）を設定した。
また、評価は主にインタビューと作業中発話に基づき定性的な知見を報告する。

\subsection{システム構成}
本システムは Streamlit を用いてWebアプリとして実装した。
ユーザーは粘土模型の写真（必要に応じてマスク画像）とテキストを入力し、システムはLLM（OpenAI GPT-4o）により入力を解釈して解釈案と属性候補を生成する。
ユーザーが選択した解釈に基づき、画像生成モデル（GPT-Image-1）によってデザイン案画像を生成し、UI上に提示する。
部分編集および合成では Inpainting API を用い、マスク領域に限定した生成や、部品と本体の接合表現を行う。生成履歴はノードとして保持し、ユーザーは過去の候補に戻って再探索できる。

「属性空間」は、形状や質感を表す属性一覧を定義した多次元空間である。
各属性は0から1の連続値で強度を表し、解釈案の提示と制約探索の双方に用いる。

\subsection{実験方法}
本実験では属性空間として形状に関する53属性を用いた。
被験者は2名（以下、被験者P、被験者Qとする）とし、日用品を対象とした2課題に取り組んでもらった。
課題は「椅子をデザインしてください」「カップをデザインしてください」であり、正解画像は設定しない。
各条件について、被験者が満足するまで生成を繰り返してよいものとし、時間制限は設けなかった。
作業の様子は録画し、実験後には半構造化インタビューを行った。
なお、被験者Pの実験は部分編集実装前のバージョンで実施したため、全体編集のみを評価対象とした。

実験条件の概要を表\ref{tab:conditions}に示す。

提案手法では、粘土画像＋テキスト入力の後、解釈案の選択と属性制約による探索を行う。
これは、粘土による形状の操作と、属性操作による詳細な制御の両方を用いる条件である。

比較Aでは、粘土画像＋テキストで初期生成した後、テキストプロンプトのみで反復生成する（Image-to-Image）。
本条件は、提案手法から「属性制約（特徴ベクトル）」を除外した構成に相当する。
直前の生成画像を次の入力として用いることで形状の維持を図るものの、言語指示のみによる微調整がいかに全体構造の崩れを招きやすく、制御性が低く探索の難易度が高いかを検証するために設定した。

比較Bでは、テキストのみを入力し、反復時も直前の生成画像を入力に含めない（Text-to-Image）。
本条件は、提案手法から「粘土による物理的形状入力」を除外した構成に相当する。
物理的な手がかりがない場合、多義的な言語のみではユーザーの意図する具体的な形状特徴を反映させることが困難であることを検証するために設定した。

\begin{table}[H]
  \centering
  \caption{実験条件の概要}
  \label{tab:conditions}
  \begin{tabular}{|p{0.18\linewidth}|p{0.36\linewidth}|p{0.36\linewidth}|}
    \hline
    条件 & 初期入力 & 反復時の入力・操作 \\
    \hline
    提案手法 & 粘土画像＋テキスト & 解釈案の選択、特徴ベクトル生成、属性制約による探索（必要に応じて部分編集／合成） \\
    \hline
    比較A & 粘土画像＋テキスト & テキストプロンプトによる反復生成（直前の生成画像を次の入力に含める） \\
    \hline
    比較B & テキストのみ & テキストプロンプトによる反復生成（直前の生成画像は入力しない） \\
    \hline
  \end{tabular}
\end{table}

実験後インタビューは半構造化形式とし、手を動かすことが与える影響（Q1）、初期出力のギャップ（Q2）、微調整の制御感（Q3）、システムからの提案の有用性（Q4）、言語化の壁（Q5）、属性リスト選択の不自由さ（Q6）、および最終成果物への反映度（Q7）を中心に質問した。インタビュー項目の概要を表\ref{tab:interview}に示す。

\begin{table}[H]
  \centering
  \caption{実験後インタビュー項目の概要}
  \label{tab:interview}
  \begin{tabular}{|p{0.16\linewidth}|p{0.76\linewidth}|}
    \hline
    質問 & ねらい \\
    \hline
    Q1 & 手を動かす（粘土造形）ことが、発想やイメージの変化に与える影響を把握する。 \\
    \hline
    Q2 & 最初の生成結果に対する「意図が伝わっていない」感覚が、どの条件で強いかを把握する。 \\
    \hline
    Q3 & 微調整における制御可能性（属性制約による操作とテキスト指示による操作の差）を把握する。 \\
    \hline
    Q4 & システムが提示する属性・解釈が、ユーザーの発想拡張や方向性決定に寄与するかを把握する。 \\
    \hline
    Q5 & テキスト指示における言語化困難（言葉が見つからず妥協する場面）を把握する。 \\
    \hline
    Q6 & 「属性リストから選ぶ」形式が制約として働く場面や、より自由な表現要求を把握する。 \\
    \hline
    Q7 & 最終成果物に対する納得感が高い条件を把握する。 \\
    \hline
  \end{tabular}
\end{table}

定性的評価として、作業中の発話と実験後インタビューを内容分析し、操作感、制御可能性、満足度、発想支援効果などの観点で整理した。

\subsection{結果の説明}
本節では、実験後インタビュー（被験者P、被験者Q）および実験中の発言記録（被験者Q）をもとに、提案手法の有用性と課題を整理する。結果は、（1）物理的形状入力、（2）解釈提案と属性制約、（3）比較手法との差異、の三つの観点からまとめる。

\subsubsection{物理的形状入力（粘土）の影響}
粘土を用いた物理的形状入力は、発想の誘発という観点で肯定的に捉えられた。
被験者Qは、テキストのみでの試行錯誤では「いや違うんやけどな」というストレスが生じやすい一方、粘土では手を動かしながら考えられるため負荷が低いと述べている。
また「新しいものができるのはやっぱりこねてるとき」という発言は、粘土操作が単なる入力手段ではなく、アイデア創出を支える活動そのものになっている可能性を示唆する。

一方で、粘土による造形は、ユーザーの手先や造形経験に依存する。
被験者Pは「頭でイメージしたことを手作業で表現することが難しい」と述べ、粘土で大体の形を伝えること自体が困難な場面があった。
ただし同時に、言葉で説明するよりも「手で作って伝えたほうが構成要素を伝えるのが楽」とも述べており、少なくとも全体構成（コンポジション）の共有には有効である可能性がある。

また、物理入力にはノイズの問題が伴う。
被験者Qの実験では、粘土の割れ目が「特徴」として解釈され生成結果に影響したという発言が複数回見られた。
これは、粘土表面のシワや継ぎ目など、ユーザーの意図しない要素が画像として入力されることで、生成モデルがそれを強い手がかりとして利用してしまうためである。
従って、撮影条件のガイドや入力画像の前処理など、物理入力を前提としたノイズ対策が重要となる。

\subsubsection{解釈提案と属性制約の影響}
提案手法の特徴である解釈提案は、探索開始の足がかりとして評価された。
被験者Qは「解釈の提示はすごい嬉しい」と述べており、抽象的な指示を具体的な観点（属性の組み合わせ）へ落とし込むことで、次に何を試せばよいかが見えやすくなると考えられる。

一方で、提示テキストや画像生成の待ち時間が長い場合には、精神的な負担やストレスとなり得る。 
被験者Qは『解釈が文字ならイメージできない』と述べており、生成結果が不明瞭な状況下でテキストを読み込む作業にはストレスが伴うことを示唆している。
また、自分が作った粘土の形をAIが意図通りに解釈しているか不安であるため、各解釈案に参照画像を添えてほしいという要望も見られたが、同時に『複数の画像生成は遅すぎるため現実的ではない』という諦めの様子も観察された。
したがって、解釈提示を支援として機能させるためには、画像生成を待たずとも直感的に判断できる表現の簡潔さや提示数の調整に加え、待ち時間をストレスに感じさせない進捗提示（プログレスバー等）の実装といった、システム性能とUXの両面からの改善が重要である。

また、属性制約による探索は、思いがけない発見の効果が観察された。
被験者Qは「頭になかった視点を示してくれた。」「ぐるぐるにしてくれたのが面白かった。」と述べ、ユーザーの想定外の候補が好意的に受け取られた例がある。
この点は、本研究が目指す「探索的デザイン支援」の有用性を支持する。

しかし制御性の観点では課題も明確であった。被験者Pは「場所の指定ができない」と述べ、デザイン全体に対する属性選択（制約）による操作では局所的な意図を伝えにくいことを指摘した。
被験者Qも「明確な意図があってコントロールしたいなら制約による探索はやってられない」と述べており、デザインを絞り込む局面においては操作が冗長に感じられる可能性がある。

さらに言語的な解釈と実際の生成結果とのズレも大きな問題として確認された。
被験者Pは「自分の想像する『くびれ』とAIが考える『くびれ』が違う」と発言し、被験者Qも「長い文章でも結局出力はズレる」と述べ、両者ともに選んだ解釈テキストから想起されるイメージと実際に生成される画像との間に大きなズレがあることを指摘した。
これを受けて被験者Qは「部分的な修正に関しては、解釈を介さずに制約（パラメータ）で直接操作したい。」という意向を示した。
これは、LLMによる言語的解釈を介するプロセスがあることで、局所的な編集においてユーザーの意図通りの制御が難しく、ストレス要因となり得ることを示唆しており、解釈を介さない直接的な部分編集機能の必要性を裏付けるものである。

\subsubsection{比較手法との差異}
比較A（粘土＋テキストのみ）は、明確な修正意図がある場合に制御しやすい傾向が見られた。
被験者Pは「比較手法Aの方が良かった」と述べ、特に場所の指定など局所的な指示を言葉で与えられる点を評価した。

比較B（テキストのみ）は、発想の拡散という観点で有用であった。
被験者Pは「ポンと出してくれるからたまにいい」と述べており、例えば「日本刀」のような具体的なモチーフを、造形の手間をかけずに短時間で出現させられる点は、探索の初期段階において有用であることが示唆された。
ただし、形状の絞り込みには課題が残り、結果として比較Aや提案手法へ移行して調整する方が良いという意見も得られた。
また、被験者からは「適切な言葉を探すのがしんどい」「諦めて妥協した」という発言が相次いだ。
視覚的なコンテキスト（粘土）がない状態でのプロンプト探索は、ユーザーに適切な言語表現を探すための思考やストレスを強いることが確認された。

\subsubsection{生成結果の質的評価}
本節では、実験中に得られた生成画像の具体例を示し、提案手法の特性を視覚的に検証する。

図\ref{fig:chair-example}に、システムによる解釈のズレがユーザーの発想支援に寄与した生成例（椅子タスク）を示す。
入力された粘土模型に対し、ユーザーは「流線型、有機的」といったクエリを与え、システムはこれに基づいて「背もたれを螺旋状にねじる」という解釈案を提示した。
しかし、実際の初期生成画像（図中）では、背もたれではなく「脚部が螺旋状になる」という、解釈テキストとは異なるデザインが出力された。
意外にも、被験者はこの「誤読」によるデザインを「掃除はしにくそうだが、面白い」と肯定的に評価した。
これは、LLMの言語的解釈と画像生成モデルの視覚的出力の間にギャップが存在することを示しているが、同時に、その予測不可能性がユーザーの想像を超えたデザイン案（セレンディピティ）をもたらす可能性も示唆している。
その後、背もたれ部分の粘土由来の不自然な凹凸を軽減するため、ユーザーは「折り曲げ（surface:bent） $\ge$ 0.85」という制約を追加した。その結果、螺旋状の脚部という特徴を維持しつつ、全体的に滑らかで洗練された最終生成物（図右）が得られた。

一方で、物理入力特有の課題が顕在化した例を図\ref{fig:mug-example}に示す。
これはマグカップタスクにおける生成例であるが、粘土表面の微細な凹凸や意図しない「シワ」が、画像生成エンジンによって「溶けたような複雑なテクスチャ」や「装飾的な文様」として過剰に解釈されている。
被験者Pが指摘した「不要なディテール」や「不気味さ」は、このような物理入力のノイズが直接的に生成結果へ反映されたことに起因すると考えられる。

\begin{figure}[H]
  \centering
  \includegraphics[width=0.32\linewidth]{/Users/umeking5/university/lab_research/truecoding/data/images/initial_20260126_165907_6be4220e.jpg}
  \hfill
  \includegraphics[width=0.32\linewidth]{/Users/umeking5/university/lab_research/truecoding/data/images/generated_20260126_172339_be3fa148.png}
%   \hfill
  \includegraphics[width=0.32\linewidth]{/Users/umeking5/university/lab_research/truecoding/data/images/generated_20260126_172522_7f4e864c.png}
  \caption{ユーザーの想定を超えた提案が行われた例}
  \label{fig:chair-example}
\end{figure}

\begin{figure}[H]
  \centering
  \includegraphics[width=0.46\linewidth]{/Users/umeking5/university/lab_research/truecoding/data/images/initial_20260126_190138_a11ae3d9.jpg}
  \hfill
  \includegraphics[width=0.46\linewidth]{/Users/umeking5/university/lab_research/truecoding/data/images/generated_20260126_190711_4562329e.png}
  \caption{物理ノイズが解釈に影響を与えた例}
  \label{fig:mug-example}
\end{figure}

\subsection{考察} 
実験結果は、粘土による物理入力が単なる形状伝達の手段にとどまらず、ユーザーの思考プロセスそのものを支援していることを示した。 
被験者Qの「こねている時に新しいものができる」という発言や、被験者Pの「言葉より手で作った方が構成を伝えるのが楽」という発言は、「形を言葉で説明する」という負荷の高い作業を粘土という物理媒体に任せることで、ユーザーがアイデアの探索に集中できていることを示唆する。 
これは、第1章で述べた「言葉の空間」と「形の空間」のギャップを埋める上で、物理入力が有効な橋渡し役となっていることを裏付けるものである。

一方で、物理入力には「ユーザーの造形スキルへの依存」と「ノイズ」という二つの課題が明らかになった。 
第一に、被験者Pが「頭でイメージしたことを手作業で表現することが難しい」と述べたように、非専門家にとっては粘土で思い通りの形状を作り出すこと自体が困難であり、不完全な入力形状が生成結果の質を低下させる要因となる点である。
第二に、マグカップの生成例（図\ref{fig:mug-example}）で見られたように、ユーザーが意図していない粘土の割れ目やシワを、AIが「模様」や「溶けたような質感」などとして過剰に解釈してしまう現象である。 
人間であれば無視できる「粗さ」をAIが特徴として拾ってしまう点は、物理とデジタルを接続する際の障壁となる。 
これに対しては、撮影時の照明ガイドや、画像処理による表面の平滑化、あるいは「これは粘土である」という前提知識をAIに強く持たせるようなプロンプト設計などの対策が必要である。

AIによる解釈と生成結果の間に生じる「ズレ」は、ユーザーの目的によって「ストレス」にも「発見」にもなり得ることが確認された。 
被験者Pのように「ここを直したい」というデザインの絞り込みのための明確な修正意図を持っている場合、意図した通りのデザインが出ないことは強いストレスとなる。 
一方で、被験者Qのように「何か面白い案が欲しい」という探索的な意図を持っている場合、椅子の脚が螺旋状になった例（図\ref{fig:chair-example}）のように、AIの「誤読」がユーザーの想像を超えたアイデアとして歓迎される場合がある。

この結果は、「ズレ」を許容できる初期の探索フェーズと、精密な制御が求められる仕上げのフェーズでは、求められる支援の形が異なることを示唆している。 
したがって、デザインの進行度に応じて、AIの解釈を挟んで大きく変化させるモードと、解釈を挟まずにパラメータや直接指示で厳密に制御するモードを、ユーザーが柔軟に切り替えられる設計が重要であると考えられる。

比較手法B（テキストのみ）での「言葉が見つからず妥協した」という結果と比較すると、提案手法における「解釈提案」は、ユーザーがゼロから言葉を探す負担を軽減する足がかりとして機能していた。 
ただし、実験で確認されたように、提示された解釈テキストが長すぎると読む気が失せたり、イメージしがたかったりする問題がある。 
ユーザーは詳細な説明よりも、直感的に判断できるキーワードや、大まかな方向性の提示を求めている傾向が見られた。 
このことから、AIからの提案は、厳密な言語定義を目指すよりも、ユーザーが「これじゃない」「こっちに近い」と感覚的に選別できるような、軽量で視覚的な対話プロセスへ改善する必要がある。




% 5章
\section{結論} 
本研究では、非専門家が抱く曖昧なデザインイメージを具体化するために、粘土による物理的な形状入力と、生成AIによる解釈・生成を組み合わせた探索的デザイン支援手法を提案した。 
提案手法は、言語だけでは伝えきれない「形状のニュアンス」を粘土で入力し、言語だけでは気づかない「言葉の多義性」をAIによる解釈提案で補うことで、ユーザーの意図を汲み取りながらデザインを具体化することを目指した。

定性的評価の結果、以下の知見が得られた。 
第一に、粘土をこねるという身体的な行為が、言葉に詰まるユーザーの思考を助け、アイデアの創出を促進することである。 
第二に、AIが提示する解釈案は、ユーザーにとって探索を始めるための有効な「きっかけ」になる一方で、生成画像との間にズレが生じると、修正局面（収束）においてはストレス要因となることである。 
第三に、粘土ならではの難しさである。粘土を用いたものづくりに慣れていないと、思いついた形を綺麗に作れず、その不格好さが生成結果の悪化につながってしまう。
また、粘土についたシワやひび割れが、AIによって勝手に「模様」や「装飾」だと勘違いされてしまい、意図しないデザインになってしまう課題も確認された。

これらの結果から、本手法は「明確なゴールに向かって形を整える作業」よりも、「ぼんやりしたイメージから可能性を広げ、思いがけないデザインに出会う作業」に特に適していると結論付けられる。

今後の展望として、以下の3点が挙げられる。 
第一に、物理入力のノイズ対策である。
画像処理による補正に加え、粘土以外の身近な日用品（ブロックや箱など）を「見立て」て入力できるような、より手軽な入力手段への拡張を検討する。 
第二に、制御性の向上である。AIの解釈を介さない直接的な部分編集機能を強化し、探索の後半でユーザーが細部を納得いくまで調整できる仕組みを整備する。 
第三に、インタフェースの改善である。AIの処理待ち時間を軽減する設計や、解釈案をテキストだけでなく簡易な画像とセットで提示するなど、直感的な判断を支えるUIの実装を進める。

本研究が提案した「物理的な直感」と「AIの表現力」を融合させるアプローチは、専門的なスキルを持たない人々が、自身の内面にあるイメージを形にし、創造的な活動を楽しむための新たな一歩になると期待する。


\acknowledgments
% ⟨ 謝辞 ⟩
%
% \bibliographystyle{kuisunsrt} % または kuissort
% \bibliography{references}
%
% 付録があれば \appendix/\Appendix に続いてここに記す。
\end{document}
